\section{Print the Longest Increasing Subsequence}
\label{sec:dp-print-lis}

\problemstatement{Given $\textit{arr}$, construct (not just measure the
length of) one actual longest increasing subsequence.}

\begin{patternbox}
Exactly as in the DP on Strings chapter's Print LCS section, reconstructing the actual
solution needs the full tabulation table kept around, plus a
\textbf{parent pointer} recorded at each index, tracking \emph{which
earlier index} extended into it -- since the table itself only stores
chain \emph{lengths}, not which specific chain produced them.
\end{patternbox}

\subsection*{Understanding the reconstruction}

Alongside $\textit{dp}[i]$ (the LIS length ending at $i$), maintain
$\textit{parent}[i]$: the index $j<i$ that $\textit{dp}[i]$ was
extended from (or $i$ itself, if $\textit{dp}[i]=1$ and nothing extends
it). After filling the table, find the index with the maximum
$\textit{dp}$ value, then follow $\textit{parent}$ pointers backward
until reaching an index whose parent is itself, collecting elements
along the way.

\subsection*{Algorithm}

\begin{enumerate}
  \item Initialize $\textit{dp}[i] \gets 1$, $\textit{parent}[i] \gets
        i$ for all $i$.
  \item For $i$ from $0$ to $n-1$: for $j$ from $0$ to $i-1$: if
        $\textit{arr}[j]<\textit{arr}[i]$ and $1+\textit{dp}[j] >
        \textit{dp}[i]$: update $\textit{dp}[i] \gets 1+\textit{dp}[j]$;
        $\textit{parent}[i] \gets j$.
  \item Find $\textit{lastIndex}$ maximizing $\textit{dp}[i]$.
  \item Walk backward from $\textit{lastIndex}$ via \textit{parent}
        until $\textit{parent}[i]=i$, collecting each $\textit{arr}[i]$
        visited; reverse the collected list.
\end{enumerate}

\subsection*{Dry run}

\dryrun{Using Section~\ref{sec:dp-lis}'s table for
$\textit{arr}=[10,9,2,5,3,7,101,18]$.}

$\textit{dp}=[1,1,1,2,2,3,4,4]$, with $\textit{parent}[3]=2$ (chain
$5$ extends $2$), $\textit{parent}[5]=3$ (chain $7$ extends $5$),
$\textit{parent}[6]=5$ (chain $101$ extends $7$). The maximum
$\textit{dp}$ value is $4$, first achieved at index $6$. Walking
backward: $6 \to 5 \to 3 \to 2$, collecting $101, 7, 5, 2$; reversing
gives $[2,5,7,101]$.

\subsection*{C++ implementation}

\begin{lstlisting}
#include <bits/stdc++.h>
using namespace std;

vector<int> printLIS(vector<int>& arr) {
    int n = arr.size();
    vector<int> dp(n, 1), parent(n);
    for (int i = 0; i < n; i++) parent[i] = i;

    int lastIndex = 0;
    for (int i = 0; i < n; i++) {
        for (int j = 0; j < i; j++) {
            if (arr[j] < arr[i] && 1 + dp[j] > dp[i]) {
                dp[i] = 1 + dp[j];
                parent[i] = j;
            }
        }
        if (dp[i] > dp[lastIndex]) lastIndex = i;
    }

    vector<int> result;
    int cur = lastIndex;
    while (parent[cur] != cur) {
        result.push_back(arr[cur]);
        cur = parent[cur];
    }
    result.push_back(arr[cur]);
    reverse(result.begin(), result.end());
    return result;
}

int main() {
    vector<int> arr = {10, 9, 2, 5, 3, 7, 101, 18};
    for (int x : printLIS(arr)) cout << x << " ";
    cout << endl; // 2 5 7 101
    return 0;
}
\end{lstlisting}

\begin{complexitybox}
\textbf{Time Complexity.} $O(n^2)$ to build the table, plus $O(n)$ for
the backward walk.
\textbf{Space Complexity.} $O(n)$ for the \textit{dp} and
\textit{parent} arrays.
\end{complexitybox}

\begin{takeawaybox}
``Print the actual chain'' problems on an ending-here table always need
a parent-pointer array recorded \emph{during} the forward fill, then a
backward walk from whichever index holds the global maximum -- the same
reconstruction idea as Print LCS, adapted to this chapter's different
table shape.
\end{takeawaybox}
